\documentclass[10pt,a4paper]{amsart}
\usepackage[margin=19mm]{geometry}
\usepackage{amsmath,amssymb,mathtools}
\usepackage[scr=rsfs,cal=euler]{mathalfa}
\usepackage{xfrac}
\usepackage[unicode,hidelinks,bookmarks=false]{hyperref}
\newcommand{\ie}{i.e.\ }
\newcommand{\cf}{cf.\ }
\newcommand{\resp}{respectively\ }
\newcommand{\Subsection}{\subsection}
\providecommand{\nobreakdash}{\nobreak\hbox{-}}
\setlength{\emergencystretch}{2em}
\multlinegap0pt
\usepackage{bm}
\usepackage{bbm}
\usepackage{dsfont}
\usepackage{xspace}
\usepackage{cancel}
\usepackage{mathtools}
\usepackage{amsthm}
\makeatletter
\@ifundefined{theorem}{\newtheorem{theorem}{Theorem}}{}
\makeatother
\title[On Kernels and Covariance Structures in Hilbert Space Gaussi]{On Kernels and Covariance Structures in Hilbert Space Gaussian Processes}
\author{Saeed Hashemi Sababe}
\date{Source: arXiv:2511.00142v1 (CC BY 4.0)}
\begin{document}
\maketitle

\noindent\emph{Source and licence.} On Kernels and Covariance Structures in Hilbert Space Gaussian Processes. Version arXiv:2511.00142v1 (CC BY 4.0), distributed under the Creative Commons Attribution 4.0 International licence, as recorded in the arXiv licence metadata for this exact version and shown on the abstract page https://arxiv.org/abs/2511.00142v1. Full rights notice: licenses/arxiv-2511.00142-CC-BY-4.0.txt.\par\smallskip

\noindent\textit{Editorial scope.} Counted unit: the Theorem whose source label is \texttt{thm:covariance\_operators} in the article (source0.tex lines 306--332, source label \texttt{thm:covariance\_operators}), delivered together with the article's own complete original proof of it (source0.tex lines 334--373, one \texttt{proof} environment of 40 lines). No mathematics is written, repaired or re-derived: statement and proof are the article's own text line for line. Required context retained uncounted: none. Every definition, display and label of those lines is retained unchanged. Omitted from this package and not redistributed: 1--305, 333--333, 374--747. Nothing inside a retained excerpt is omitted or altered: every retained body is delivered exactly as the article's lines are written, so no omission receipt of any kind accompanies this package. Citations: the retained excerpts carry no citation command at all, so no bibliography entry of the article is transcribed and no reference is redistributed. Reference adaptation: none was needed and none was made; every reference inside the delivered unit resolves inside it, so no unresolved reference or citation is produced. Printed slips kept literal: no printed word, symbol, hypothesis or number of the counted statement or of its proof was altered. Layout and class: the article's own class is \texttt{amsart} with packages including amsmath, amssymb, graphicx, pgfplots, caption, cite, cases, booktabs, hyperref; the wrapper is the lane's pinned \texttt{amsart} editorial wrapper at 10pt on A4, which restates the theorem-like environments the retained text needs and adds the article's own macro definitions that the retained text needs (none), with extra wrapper packages (bm, bbm, dsfont, xspace, cancel, mathtools). The delivered numbering is the unit's own. Assets: no retained span contains a figure environment, a graphics-inclusion command, a PDF-page inclusion, a table or a TikZ picture; no image file is copied into this package, the package declares no asset dependency and no image file is redistributed with it. Licence: version arXiv:2511.00142v1 of this article is distributed under the Creative Commons Attribution 4.0 International licence, as recorded in the arXiv licence metadata for this exact version and shown on the abstract page https://arxiv.org/abs/2511.00142v1. A full rights notice is delivered at licenses/arxiv-2511.00142-CC-BY-4.0.txt. Characters: every retained excerpt is pure ASCII, so the delivered engine, XeLaTeX with its default Latin Modern text font, reports no missing character.
\begin{theorem}
\label{thm:covariance_operators}
Let \( \{ V_s \}_{s \in S} \) be the family of operators \( V_s: H \to H_{\tilde{K}} \) defined by
\[
V_s a = \tilde{K}(\cdot, (s, a)): X \to \mathbb{C}, \quad a \in H.
\]
Then the following properties hold:
\begin{enumerate}
    \item The covariance operator \( \Sigma_s: H \to H \) is self-adjoint and positive semi-definite for all \( s \in S \).
    \item The operator \( V_s V_s^*: H_{\tilde{K}} \to H_{\tilde{K}} \) acts as the covariance operator in the RKHS \( H_{\tilde{K}} \), and it satisfies:
    \[
    V_s V_s^* \tilde{K}(\cdot, (t, b)) = \tilde{K}(\cdot, (s, K(s,t) b)).
    \]
    \item The covariance operator \( \Sigma_s \) induces a bounded operator on \( H_{\tilde{K}} \), and for all \( s, t \in S \) and \( a, b \in H \),
    \[
    \langle \Sigma_s a, b \rangle_H = \langle V_s a, V_s b \rangle_{H_{\tilde{K}}}.
    \]
    \item The covariance operator \( \Sigma_s \) can be expressed as:
    \[
    \Sigma_s = V_s^* V_s.
    \]
    \item The operators \( V_s \) are isometric if \( K(s,s) = I_H \), and the covariance operators satisfy:
    \[
    \langle \Sigma_s a, a \rangle_H = \| V_s a \|_{H_{\tilde{K}}}^2.
    \]
\end{enumerate}
\end{theorem}

\begin{proof}
We prove the items as the following:
\subsubsection*{1. Self-adjointness and Positive Semi-Definiteness:} The operator \( \Sigma_s: H \to H \), defined by \( \Sigma_s a = K(s,s)a \), is clearly self-adjoint because \( K(s,s) \) is self-adjoint, i.e.,
    \[
    \langle \Sigma_s a, b \rangle_H = \langle K(s,s)a, b \rangle_H = \langle a, K(s,s)b \rangle_H = \langle a, \Sigma_s b \rangle_H.
    \]
    Moreover, for any \( a \in H \),
    \[
    \langle \Sigma_s a, a \rangle_H = \langle K(s,s) a, a \rangle_H \geq 0,
    \]
    showing that \( \Sigma_s \) is positive semi-definite.
    
\subsubsection*{2. Action of \( V_s V_s^* \) as a Covariance Operator:} The operator \( V_s V_s^* \) maps vectors in \( H_{\tilde{K}} \) to vectors in \( H_{\tilde{K}} \). By definition of \( V_s \) and \( V_s^* \),
    \[
    V_s V_s^* \tilde{K}(\cdot, (t, b)) = V_s (K(s, t) b) = \tilde{K}(\cdot, (s, K(s,t) b)).
    \]
    This shows that \( V_s V_s^* \) acts as a covariance operator in \( H_{\tilde{K}} \).
    
\subsubsection*{3. Induced Covariance Operator in \( H_{\tilde{K}} \):} We have
    \[
    \langle V_s a, V_s b \rangle_{H_{\tilde{K}}} = \langle \tilde{K}(\cdot, (s, a)), \tilde{K}(\cdot, (s, b)) \rangle_{H_{\tilde{K}}} = \langle a, K(s,s) b \rangle_H = \langle \Sigma_s a, b \rangle_H,
    \]
    so the covariance operator \( \Sigma_s \) induces a bounded operator on \( H_{\tilde{K}} \).
    
\subsubsection*{4. Expression of \( \Sigma_s \) as \( V_s^* V_s \):} From the above, we have
    \[
    \langle \Sigma_s a, b \rangle_H = \langle V_s a, V_s b \rangle_{H_{\tilde{K}}},
    \]
    which implies that \( \Sigma_s = V_s^* V_s \), since both sides act as the covariance operator on \( H \).
    
\subsubsection*{5. Isometry and Covariance Operator Norm:} If \( K(s,s) = I_H \), then for all \( a \in H \),
    \[
    \langle \Sigma_s a, a \rangle_H = \langle a, a \rangle_H,
    \]
    so \( V_s \) is an isometry. Furthermore, we have
    \[
    \| V_s a \|_{H_{\tilde{K}}}^2 = \langle \Sigma_s a, a \rangle_H.
    \]

\end{proof}

\end{document}
